\section{Proofs of Implications}
\label{sec:impls-extra}

\subsection{Among Derived Notions}
\label{sec:impls-extra:among-derived}

\begin{remark}
\label{thm:impl:epistemic}
For any fairness notion $F$, if an allocation is $F$-fair to an agent $i$,
then it is also epistemic-$F$-fair to agent $i$.
If an allocation is epistemic-$F$-fair to an agent $i$,
then it is also minimum-$F$-share-fair to agent $i$.
If there are only two agents, then an allocation is epistemic-$F$-fair to an agent $i$
iff it is $F$-fair to agent $i$.
\end{remark}

\begin{remark}
\label{thm:impl:groupwise}
For any fairness notion $F$, if an allocation is groupwise-$F$-fair to an agent $i$,
then it is also pairwise-$F$-fair to agent $i$ and $F$-fair to agent $i$.
When there are only two agents, all three of these notions are equivalent.
\end{remark}

\begin{lemma}
\label{thm:impl:ext-to-epistemic}
Let $\Omega$ be a set containing pairs of the form $(\Ical, A)$,
where $\Ical$ is a fair division instance and $A$ is an allocation for $\Ical$.
For any two fairness notions $F_1$ and $F_2$, let $\phi(F_1, F_2)$ be the proposition
``\,$\forall (\Ical, A) \in \Omega$, for every agent $i$ in $\Ical$,
$A$ is $F_2$-fair to $i$ whenever $A$ is $F_1$-fair to $i$". Then
$\phi(F_1, F_2) \implies \phi(\textrm{epistemic-}F_1, \textrm{epistemic-}F_2)
\textrm{ and } \phi(\textrm{min-}F_1\textrm{-share}, \textrm{min-}F_2\textrm{-share})$.
\end{lemma}
\begin{proof}
Suppose $\phi(F_1, F_2)$.
Pick any $(\Ical, A) \in \Omega$.
Let $\Ical \defeq ([n], [m], (v_i)_{i=1}^n, w)$.
Pick any $i \in [n]$.

Suppose $A$ is epistemic-$F_1$-fair to $i$.
Let $B$ be $i$'s epistemic-$F_1$-certificate for $A$.
Then $B$ is $F_1$-fair to $i$ and $A_i = B_i$.
By $\phi(F_1, F_2)$, $B$ is also $F_2$-fair to $i$.
Hence, $B$ is $i$'s epistemic-$F_2$-certificate for $A$.
Therefore, $\phi(\textrm{epistemic-}F_1, \textrm{epistemic-}F_2)$ holds.

Suppose $A$ is min-$F_1$-share-fair to $i$.
Let $B$ be $i$'s min-$F_1$-share-certificate for $A$.
Then $B$ is $F_1$-fair to $i$ and $v_i(A_i) \ge v_i(B_i)$.
By $\phi(F_1, F_2)$, $B$ is also $F_2$-fair to $i$.
Hence, $B$ is $i$'s min-$F_2$-share-certificate for $A$.
Therefore, $\phi(\textrm{min-}F_1\textrm{-share}, \textrm{min-}F_2\textrm{-share})$ holds.
\end{proof}

\subsection{Among EF, EFX, EF1}
\label{sec:impls-extra:among-ef-efx-ef1}

Here we look at implications among EF, EFX, EF1, and their epistemic variants.

\begin{remark}[EF $\fimplies$ EFX+EF1]
\label{thm:impl:ef-to-efx+ef1}
If an allocation is EF-fair to agent $i$, then it is also EFX-fair to $i$ and EF1-fair to $i$.
\end{remark}

Because of how we define EFX (\cref{defn:efx}),
it's not always true that EFX implies EF1.
However, it's true for many common settings, as the following lemma shows.

\begin{lemma}[EFX $\fimplies$ EF1]
\label{thm:impl:efx-to-ef1}
For the fair division instance $([n], [m], (v_i)_{i=1}^n, w)$,
let $A$ be an allocation where agent $i$ is EFX-satisfied.
Then agent $i$ is EF1-satisfied in these scenarios:
\begin{tightenum}
\item $v_i$ is additive.
\item $v_i$ is doubly strictly monotone, i.e., $[m] = G \cup C$, $v_i(g \mid \cdot) > 0$
    for every $g \in G$, and $v_i(c \mid \cdot) < 0$ for every $c \in C$.
\item Agents have equal entitlements, $v_i$ is submodular, and all items are goods for agent $i$.
\item $v_i$ is submodular and all items are chores for agent $i$.
\end{tightenum}
\end{lemma}
\begin{proof}
Suppose $i$ is EFX-satisfied but not EF1-satisfied.
Suppose $i$ EF1-envies $j$.

Since $i$ EF1-envies $j$, we get that for all $t \in A_j$, we have
\[ \frac{v_i(A_i)}{w_i} < \frac{v_i(A_j \setminus \{t\})}{w_j}. \]
Since $i$ is EFX-satisfied, we get that
for all $t \in A_j$ such that $v_i(t \mid A_i) > 0$, we have
\[ \frac{v_i(A_i)}{w_i} \ge \frac{v_i(A_j \setminus \{t\})}{w_j}. \]
Hence, for all $t \in A_j$, we get $v_i(t \mid A_i) \le 0$.

Since $i$ EF1-envies $j$, we get that for all $t \in A_i$, we have
\[ \frac{v_i(A_i \setminus \{t\})}{w_i} < \frac{v_i(A_j)}{w_j}. \]
Since $i$ is EFX-satisfied, we get that
for all $t \in A_i$ such that $v_i(t \mid A_i \setminus \{t\}) < 0$, we have
\[ \frac{v_i(A_i \setminus \{t\})}{w_i} \ge \frac{v_i(A_j)}{w_j}. \]
Hence, for all $t \in A_i$, we get $v_i(t \mid A_i \setminus \{t\}) \ge 0$.

If $v_i$ is additive, we get $v_i(A_i) \ge 0 \ge v_i(A_j)$, which is a contradiction.
%
If $v_i$ is doubly strictly monotone, then
all items in $A_j$ are chores and all items in $A_i$ are goods.
Hence, $v_i(A_i) \ge 0 \ge v_i(A_j)$, which is a contradiction.

Suppose all agents have equal entitlements, all items are goods for $i$, and $v_i$ is submodular.
Let $A_j \defeq \{g_1, \ldots, g_k\}$. Then
\[ v_i(A_j \mid A_i) = \sum_{t=1}^k v_i(g_t \mid A_i \cup \{g_1, \ldots, g_{t-1}\})
\le \sum_{t=1}^k v_i(g_t \mid A_i) \le 0. \]
Hence, $v_i(A_j) \le v_i(A_i \cup A_j) = v_i(A_i) + v_i(A_j \mid A_i) \le v_i(A_i)$.
This is a contradiction.

Suppose all items are chores for $i$ and $v_i$ is submodular.
Let $A_i = \{c_1, \ldots, c_k\}$. Then
\[ v_i(A_i) = \sum_{t=1}^k v_i(c_t \mid \{c_1, \ldots, c_{t-1}\})
\ge \sum_{t=1}^k v_i(c_t \mid A_i \setminus \{c_t\}) \ge 0. \]
Hence, $v_i(A_i) \ge 0 \ge v_i(A_j)$, which is a contradiction.

A contradiction implies that it's impossible for agent $i$ to be
EFX-satisfied but not EF1-satisfied.
\end{proof}

\begin{lemma}[MXS $\fimplies$ EF1 for $n=2$]
\label{thm:impl:mxs-to-ef1-n2}
Let $([2], [m], (v_1, v_2), w)$ be a fair division instance with indivisible items.
If $v_1$ is additive and agent 1 is MXS-satisfied by allocation $A$,
then agent 1 is also EF1-satisfied by $A$.
\end{lemma}
\begin{proof}
Suppose $A$ is MXS-fair to agent 1 but not EF1-fair to her.
Then agent 1 envies agent 2 in $A$, so $v_1(A_1) < v_1(A_2)$.
Let $B$ be agent 1's MXS-certificate for $A$. Then $v_1(B_1) \le v_1(A_1)$.
Moreover, $v_1(A_2) = v_1([m]) - v_1(A_1) \le v_1([m]) - v_1(B_1) = v_1(B_2)$.
Hence, we get $v_1(B_1) \le v_1(A_1) < v_1(A_2) \le v_1(B_2)$.

Let $G \defeq \{g \in [m]: v_1(g) > 0\}$ and $C \defeq \{c \in [m]: v_1(c) < 0\}$.
Let $\max(\emptyset) \defeq -\infty$ and $\min(\emptyset) \defeq \infty$.

Since agent 1 is EFX-satisfied by $B$ and not EF1-satisfied by $A$,
for every $\ghat \in A_2$, we get
\begin{align*}
& \frac{v_1(A_2) - v_1(\ghat)}{w_2} > \frac{v_1(A_1)}{w_1}
\ge \frac{v_1(B_1)}{w_1}
\\ &\ge \frac{1}{w_2}\left(v_1(B_2) - \min_{g \in B_2 \cap G} v_1(g)\right)
\\ &\ge \frac{1}{w_2}\left(v_1(A_2) - \min_{g \in B_2 \cap G} v_1(g)\right).
\end{align*}
Hence, for every $\ghat \in A_2$, we get $v_1(\ghat) < \min_{g \in B_2 \cap G} v_1(g)$.
Hence, $A_2 \cap G$ and $B_2 \cap G$ are disjoint, so $A_2 \cap G \subseteq B_1 \cap G$.
Let $d_i \defeq -v_i$ for all $i$. Then for every $\chat \in A_1$, we get
\begin{align*}
& \frac{d_1(A_1) - d_1(\chat)}{w_1} > \frac{d_1(A_2)}{w_2}
\ge \frac{d_2(B_2)}{w_2}
\\ &\ge \frac{1}{w_1}\left(d_1(B_1) - \min_{c \in B_1 \cap C} d_1(c)\right)
\\ &\ge \frac{1}{w_1}\left(d_1(A_1) - \min_{c \in B_1 \cap C} d_1(c)\right).
\end{align*}
Hence, for every $\chat \in A_1$, we have $d_1(\chat) < \min_{c \in B_1 \cap C} d_1(c)$.
Hence, $A_1 \cap C$ and $B_1 \cap C$ are disjoint, so $B_1 \cap C \subseteq A_2 \cap C$.
Hence,
\begin{align*}
v_1(A_2) &= v_1(A_2 \cap G) - d_1(A_2 \cap C)
\\ &\le v_1(B_1 \cap G) - d_1(B_1 \cap C) = v_1(B_1),
\end{align*}
which is a contradiction.
Hence, it's not possible for $A$ to be MXS-fair to agent 1 but not EF1-fair to her.
\end{proof}

\subsection{Among PROP-Based Notions}
\label{sec:impls-extra:among-prop-based}

\begin{lemma}[PROPx $\fimplies$ PROPavg]
\label{thm:impl:propx-to-propavg}
In a fair division instance $([n], [m], (v_i)_{i=1}^n, w)$,
if an allocation is PROPx-fair to agent $i$, then it is also PROPavg-fair to agent $i$.
\end{lemma}
\begin{proof}
Assume (for the sake of contradiction) that there is an allocation $A$ where
agent $i$ is PROPx-satisfied but not PROPavg-satisfied.
%
Since $i$ is not PROPavg-satisfied, we get $v_i(A_i) \le w_iv_i([m])$.
Since $i$ is PROPx-satisfied, we get

\begin{itemize}
\item $v_i(A_i \setminus S) > w_iv_i([m])$ for all $S \subseteq A_i$ such that
    $v_i(S \mid A_i \setminus S) < 0$.
\item $v_i(A_i \cup S) > w_iv_i([m])$ for all $S \subseteq [m] \setminus A_i$
    such that $v_i(S \mid A_i) > 0$.
\end{itemize}

Since $i$ is not PROPavg-satisfied, we get that $T \neq \emptyset$
and $v_i(A_i) + \Sum(T)/|T| \le w_iv_i([m])$ ($T$ is defined in \cref{defn:propavg}).
For each $\tau_j \in T$, we have $0 < \tau_j = v_i(S_j \mid A_i)$ for some $S_j \subseteq A_j$.
Since $A$ is PROPx-fair to $i$, we have $v_i(A_i) + \tau_j > w_iv_i([m])$ for all $\tau_j \in T$.
On averaging these inequalities, we get $v_i(A_i) + \Sum(T)/|T| > w_iv_i([m])$,
which implies that $A$ is PROPavg-satisfied, which is a contradiction.
Hence, if $i$ is PROPx-satisfied by $A$, then she is also PROPavg-satisfied by $A$.
\end{proof}

\begin{lemma}
\label{thm:submod-marginal-is-submod}
For any $X \subseteq \Omega$ and any submodular function $f: 2^{\Omega} \to \mathbb{R}$,
$f(\cdot \mid X)$ is submodular.
\end{lemma}
\begin{proof}
Let $P, Q \subseteq \Omega \setminus X$. Let $g(Y) \defeq f(Y \mid X)$. Then
\begin{align*}
& g(P) + g(Q)
\\ &= f(P \cup X) + f(Q \cup X) - 2f(X)
\\ &\ge f((P \cup X) \cup (Q \cup X)) + f((P \cup X) \cap (Q \cap X)) - 2f(X)
    \tag{by $f$'s submodularity}
\\ &= f((P \cup Q) \cup X) + f((P \cap Q) \cup X) - 2f(X)
    \tag{by De Morgan's law}
\\ &= g(P \cup Q) + g(P \cap Q).
\end{align*}
Hence, $g$ is submodular.
\end{proof}

\begin{lemma}[PROPm $\fimplies$ PROP1]
\label{thm:impl:propm-to-prop1}
For a fair division instance $([n], [m], (v_i)_{i=1}^n, w)$,
if an allocation $A$ is PROPm-fair to agent $i$, then it is also PROP1-fair to agent $i$
if at least one of these conditions holds:
\begin{tightenum}
\item $v_i$ is submodular.
\item $v_i$ is doubly strictly monotone, i.e., $[m] = G \cup C$, $v_i(g \mid \cdot) > 0$
    for every $g \in G$, and $v_i(c \mid \cdot) < 0$ for every $c \in C$.
\end{tightenum}
\end{lemma}
\begin{proof}
Suppose allocation $A$ is PROPm-fair to $i$ but not PROP1-fair to $i$. Then
\begin{tightenum}
\item\label{item:impl:propm-to-prop1:unprop}$v_i(A_i) < w_iv_i([m])$ (by PROP1 unfairness).
\item\label{item:impl:propm-to-prop1:unprop1-chores}$v_i(A_i \setminus \{c\}) \le w_iv_i([m])$
    for all $c \in A_i$ (by PROP1 unfairness).
\item\label{item:impl:propm-to-prop1:unprop1-goods}$v_i(A_i \cup \{g\}) \le w_iv_i([m])$
    for all $g \in [m] \setminus A_i$ (by PROP1 unfairness).
\item\label{item:impl:propm-to-prop1:propm-chores}$v_i(A_i \setminus \{c\}) > w_iv_i([m])$
    for all $c \in A_i$ such that $v_i(c \mid A_i \setminus \{c\}) < 0$ (by PROPm fairness).
\item\label{item:impl:propm-to-prop1:propm-goods}$T = \emptyset$ or $v_i(A_i) + \max(T) > w_iv_i([m])$
    (by PROPm fairness; c.f.~\cref{defn:propm} for the definition of $T$).
\end{tightenum}

By \ref{item:impl:propm-to-prop1:unprop1-chores} and \ref{item:impl:propm-to-prop1:propm-chores},
we get $v_i(c \mid A_i \setminus \{c\}) \ge 0$ for all $c \in A_i$.
We now show that $v_i(A_i) \ge 0$.
If $v_i$ is doubly strictly monotone, then $A_i$ only contains goods, so $v_i(A_i) \ge 0$.
Now suppose $v_i$ is submodular. Let $A_i = \{g_1, \ldots, g_k\}$. Then
\[ v_i(A_i) = \sum_{j=1}^k v_i(g_j \mid \{g_1, \ldots, g_{j-1}\})
    \ge v_i(g_j \mid A_i \setminus \{g_j\}) \ge 0. \]

Suppose $T = \emptyset$. Then $\tau_j = 0$ for all $j \in [n] \setminus \{i\}$.
Hence, for all $j \in [n] \setminus \{i\}$, we have $v_i(A_j \mid A_i) \le 0$.
If $v_i$ is doubly strictly monotone, then $[m] \setminus A_i$ contains only chores,
so $v_i([m] \setminus A_i \mid A_i) \le 0$. If $v_i$ is submodular, then
$v_i(\cdot \mid A_i)$ is subadditive by \cref{thm:submod-marginal-is-submod}, so
\[ v_i([m] \setminus A_i \mid A_i)
    \le \sum_{j \in [n] \setminus \{i\}} v_i(A_j \mid A_i) \le 0. \]
Hence, $v_i(A_i) \ge v_i([m])$.
If $v_i([m]) \le 0$, then $v_i(A_i) \ge 0 \ge w_iv_i([m])$,
and if $v_i([m]) \ge 0$, then $v_i(A_i) \ge v_i([m]) \ge w_iv_i([m])$.
This contradicts \ref{item:impl:propm-to-prop1:unprop}, so $T \neq \emptyset$.

Let $\max(T) = \tau_{\jhat} > 0$. By definition of $\tau_{\jhat}$, we get
\begin{align*}
0 < \tau_{\jhat} &= \min(\{v_i(S \mid A_i) \mid S \subseteq A_{\jhat} \textrm{ and } v_i(S \mid A_i) > 0\})
\\ &\le \min(\{v_i(g \mid A_i) \mid g \in A_{\jhat} \textrm{ and } v_i(g \mid A_i) > 0\}).
\end{align*}

\textbf{Case 1}: $v_i(g \mid A_i) \le 0$ for all $g \in A_{\jhat}$.
\\ If $v_i$ is doubly strictly monotone, then $A_{\jhat}$ only has chores,
and so $v_i(S \mid A_i) \le 0$ for all $S \subseteq A_{\jhat}$.
This contradicts the fact that $\tau_{\jhat} > 0$.
Now let $v_i$ be submodular.
Since $v_i(\cdot \mid A_i)$ is subadditive by \cref{thm:submod-marginal-is-submod},
for any $S \subseteq A_{\jhat}$, we get
$v_i(S \mid A_i) \le \sum_{c \in S} v_i(c \mid A_i) \le 0$.
This contradicts the fact that $\tau_{\jhat} > 0$.

\textbf{Case 2}: $v_i(\ghat \mid A_i) > 0$ for some $\ghat \in A_{\jhat}$.
\\ Then $\max(T) = \tau_{\jhat} \le v_i(\ghat \mid A_i)$.
By \ref{item:impl:propm-to-prop1:propm-goods}, we get
$w_iv_i([m]) < v_i(A_i) + \max(T) \le v_i(A_i \cup \{\ghat\})$.
But this contradicts \ref{item:impl:propm-to-prop1:unprop1-goods}.

Hence, it cannot happen that $i$ is PROPm-satisfied by $A$ but not PROP1-satisfied.
\end{proof}

\subsection{EF vs PROP}
\label{sec:impls-extra:ef-vs-prop}

Here we look at implications between EF (and its epistemic variants) and PROP.

\begin{lemma}[MEFS $\fimplies$ PROP, \cite{bouveret2016characterizing}]
\label{thm:impl:mefs-to-prop}
For a fair division instance $([n], [m], (v_i)_{i=1}^n, w)$,
if $v_i$ is subadditive and an allocation $A$ is MEFS-fair to $i$, then $A$ is also PROP-fair to $i$.
\end{lemma}
\begin{proof}
Let $B$ be agent $i$'s MEFS-certificate for $A$.
Then for all $j \in [n]$, we have $v_i(B_i)/w_i \ge v_i(B_j)/w_j$.
Sum these inequalities over all $j \in [n]$, weighting each by $w_j$,
to get $v_i(B_i)/w_i \ge \sum_{j=1}^n v_i(B_j)$.
Since $v_i$ is subadditive, we get $v_i([m]) \le \sum_{j=1}^n v_i(B_j)$.
Hence,
\[ \frac{v_i(A_i)}{w_i} \ge \frac{v_i(B_i)}{w_i} \ge \sum_{j=1}^n v_i(B_j) \ge v_i([m]). \]
\end{proof}

\begin{lemma}[EF $\fimplies$ GPROP]
\label{thm:impl:ef-to-gprop}
For a fair division instance $([n], [m], (v_i)_{i=1}^n, w)$,
if $v_i$ is subadditive and agent $i$ is envy-free in $A$,
then $A$ is groupwise-PROP-fair to $i$.
\end{lemma}
\begin{proof}
Let $S$ be a subset of agents containing $i$.
Let $\Ahat$ be the allocation obtained by restricting $A$ to $S$ (c.f.~\cref{defn:restricting}).
Then $i$ is also envy-free in $\Ahat$.
$i$ is also MEFS-satisfied by $\Ahat$, since $\Ahat$ is its own MEFS-certificate for agent $i$.
By \cref{thm:impl:mefs-to-prop}, agent $i$ is PROP-satisfied by $\Ahat$.
Since this is true for all $S$ containing $i$,
we get that $A$ is groupwise-PROP-fair to agent $i$.
\end{proof}

\begin{lemma}[PROP $\fimplies$ EF for idval]
\label{thm:impl:prop-to-ef-superadd-id}
In a fair division instance $([n], [m], (v_i)_{i=1}^n, w)$ with identical superadditive valuations,
a PROP allocation is also an EF allocation.
\end{lemma}
\begin{proof}
Let $v$ be the common valuation function. Let $A$ be a PROP allocation.
Then for each agent $i$, we have $v(A_i) \ge w_iv([m])$.
Suppose $v(A_k) > w_kv([m])$ for some agent $k$.
Sum these inequalities to get $\sum_{i=1}^n v(A_i) > v([m])$.
This contradicts superadditivity of $v$, so $v(A_i) = w_iv([m])$ for each agent $i$.
Hence, $v(A_i)/w_i = v([m])$ for all $i$, so $A$ is EF.
\end{proof}

\begin{lemma}[PROP $\fimplies$ EF for $n=2$]
\label{thm:impl:prop-to-ef-n2}
In the fair division instance $([2], [m], (v_1, v_2), w)$, for some agent $i$,
if $v_i$ is superadditive and agent $i$ is PROP-satisfied by allocation $A$,
then she is also envy-free in $A$.
\end{lemma}
\begin{proof}
Assume $i=1$ \wLoG. Then
\begin{align*}
& \frac{v_1(A_1)}{w_1} \ge v_1([m]) \ge v_1(A_1) + v_1(A_2)
\\ &\implies v_1(A_2) \le v_1(A_1)\left(\frac{1}{w_1} - 1\right) = w_2 \frac{v_1(A_1)}{w_1}.
\end{align*}
Hence, agent 1 does not envy agent 2.
\end{proof}

\subsection{EFX, EF1 vs PROPx, PROPm, PROP1}
\label{sec:impls-extra:efx-ef1-vs-propx-propm-prop1}

\begin{lemma}[MXS $\fimplies$ PROP1, Theorem 3 of \cite{caragiannis2023new}]
\label{thm:impl:mxs-to-prop1}
In a fair division instance with equal entitlements,
if $v_i$ is additive for some agent $i$, $v_i(g) \ge 0$ for every item $g$,
and an allocation $A$ is MXS-fair to agent $i$, then $A$ is also PROP1-fair to agent $i$.
\end{lemma}

\begin{lemma}[EEF1 $\fimplies$ PROP1, Proposition 2 of \cite{aziz2021fair}]
\label{thm:impl:eef1-to-prop1}
For a fair division instance $([n], [m], (v_i)_{i=1}^n, w)$,
if an allocation $X$ is epistemic-EF1-fair to agent $i$, then $X$ is also PROP1-fair to $i$
if one of these conditions hold:
\begin{tightenum}
\item $v_i$ is subadditive and the items are chores to agent $i$.
\item $v_i$ is additive and $w_i \le w_j$ for all $j \in [n] \setminus \{i\}$.
\item $v_i$ is additive and $n=2$.
\end{tightenum}
\end{lemma}
\begin{proof}
Suppose agent $i$ is epistemic-EF1-satisfied but not PROP1-satisfied by allocation $X$.
Let $A$ be agent $i$'s epistemic-EF1-certificate for $X$.
Then $A$ is EF1-fair to agent $i$.
If $A$ is PROP1-fair to $i$, then $X$ would also be PROP1-fair to agent $i$,
which is a contradiction. Hence, $A$ is not PROP1-fair to $i$.
Therefore, all of these hold:
\begin{tightenum}
\item $v_i(A_i) < w_iv_i([m])$.
\item $v_i(A_i \cup \{g\}) \le w_iv_i([m])$ for all $g \in [m] \setminus A_i$.
\item $v_i(A_i \setminus \{c\}) \le w_iv_i([m])$ for all $c \in A_i$.
\end{tightenum}

Since $v_i$ is subadditive, there exists $j \in [n] \setminus \{i\}$ such that $v_i(A_j) > w_jv_i([m])$
(otherwise $v_i([m]) \le \sum_{j=1}^n v_i(A_j) < \sum_{j=1}^n w_jv_j([m]) = v_i([m])$). Hence,
\[ \frac{v_i(A_i)}{w_i} < v_i([m]) < \frac{v_i(A_j)}{w_j}. \]
Hence, $i$ envies $j$. But $i$ is EF1-satisfied. Hence,
\\ $\exists c \in A_i$ such that
    $\displaystyle \frac{v_i(A_i \setminus \{c\})}{w_i} \ge \frac{v_i(A_j)}{w_j}$,
\\ or $\exists g \in A_j$ such that
    $\displaystyle \frac{v_i(A_i)}{w_i} \ge \frac{v_i(A_j \setminus \{g\})}{w_j}$.

\textbf{Case 1}: $\exists c \in A_i$ such that $v_i(A_i \setminus \{c\})/w_i \ge v_i(A_j)/w_j$.

Since $i$ is PROP1-unsatisfied and $v_i(A_j) > w_jv_i([m])$, we get
\[ \frac{v_i(A_j)}{w_j} \le \frac{v_i(A_i \setminus \{c\})}{w_i} \le v_i([m]) < \frac{v_i(A_j)}{w_j}, \]
which is a contradiction.
Hence, it's impossible for agent $i$ to be EF1-satisfied but not PROP1-satisfied by $A$ for this case.

\textbf{Case 2}: $\exists g \in A_j$ such that $v_i(A_i)/w_i \ge v_i(A_j \setminus \{g\})/w_j$.

First, we show that this case doesn't occur if all items in $A_j$ are chores.
Since $i$ envies $j$ but not EF1-envies $j$, we get
\begin{equation}
\label{eq:impl:eef1-to-prop1:gpos}
\frac{v_i(g \mid A_j \setminus \{g\})}{w_j}
= \frac{v_i(A_j) - v_i(A_j \setminus \{g\})}{w_j}
\ge \frac{v_i(A_j)}{w_j} - \frac{v_i(A_i)}{w_i} > 0.
\end{equation}
Hence, this case doesn't occur if all items in $A_j$ are chores.

Since $i$ is not PROP1-satisfied, we get
\begin{align*}
& \frac{v_i(A_i) + v_i(g \mid A_i)}{w_i} \le w_i([m]) < \frac{v_i(A_j)}{w_j}
\\ &\implies \frac{v_i(g \mid A_i)}{w_i} < \frac{v_i(A_j)}{w_j} - \frac{v_i(A_i)}{w_i}
    \le \frac{v_i(g \mid A_j \setminus \{g\})}{w_j}.
\end{align*}

Let $v_i$ be additive. Then we get
\[ \frac{v_i(g)}{w_i} < \frac{v_i(g)}{w_j} \implies w_j < w_i. \]
If $w_i \le w_j$ for all $j \in [n] \setminus \{i\}$, we get a contradiction.

Now suppose $v_i$ is additive and $n=2$.
Then the agents are $i$ and $j$. Note that $w_i + w_j = 1$.
Let $\ghat \in \argmax_{t \in A_j} v_i(t)$.
Since $i$ is not PROP1-satisfied, we get
\begin{align*}
& v_i(A_i \cup \{\ghat\}) \le w_iv_i([m]) = w_iv_i(A_i \cup \{\ghat\}) + w_iv_i(A_j \setminus \{\ghat\})
\\ &\implies \frac{v_i(A_i \cup \{\ghat\})}{w_i} \le \frac{v_i(A_j \setminus \{\ghat\})}{w_j}.
\end{align*}
Based on the assumption of Case 2, we get
\[ \frac{v_i(A_i)}{w_i} \ge \frac{v_i(A_j \setminus \{g\})}{w_j} \ge \frac{v_i(A_j \setminus \{\ghat\})}{w_j}. \]
Hence,
\[ v_i(A_i \cup \{\ghat\}) \le \frac{w_i}{w_j}v_i(A_j \setminus \{\ghat\}) \le v_i(A_i). \]
Therefore, $v_i(\ghat) \le 0$.
Since $\ghat$ is the most-valuable item in $A_j$ to $i$, we get that $v_i(g) \le 0$,
which contradicts \eqref{eq:impl:eef1-to-prop1:gpos}.
Hence, it's impossible for agent $i$ to be EF1-satisfied but not PROP1-satisfied by $A$ for this case.
\end{proof}

\begin{lemma}[EEFX $\fimplies$ PROPx, Lemma 2.1 of \cite{li2022almost}]
\label{thm:impl:eefx-to-propx}
Consider a fair division instance $([n], [m], (v_i)_{i=1}^n, w)$ where
the items are chores to agent $i$ and $v_i$ is subadditive.
If an allocation $A$ is epistemic-EFX-fair to agent $i$,
then it is also PROPx-fair to agent $i$.
\end{lemma}
\begin{proof}
Suppose $A$ is epistemic-EFX-fair to agent $i$ but not PROPx-fair to $i$.
Let $B$ be agent $i$'s epistemic-EFX-certificate for $A$.
Then $B$ is also PROPx-unfair to $i$.

Since $i$ is EFX-satisfied with $B$, we get that for all $j \in [n] \setminus \{i\}$,
\[ \min_{S \in \Scal} \frac{v_i(B_i \setminus S)}{w_i} \ge \frac{v_i(B_j)}{w_j}, \]
where $\Scal \defeq \{S \subseteq B_i: v_i(S \mid B_i \setminus S) < 0\}$.
Add these inequalities for all $j$, weighting each by $w_j$, to get
\[ \min_{S \in \Scal} \frac{v_i(B_i \setminus \{c\})}{w_i} > \sum_{j=1}^n v_i(B_j) \ge v_i([m]), \]
which implies that $B$ is PROPx-fair, a contradiction.
Hence, there can't be an allocation $A$ that is epistemic-EFX-fair to $i$
but not PROPx-fair to $i$.
\end{proof}

\begin{lemma}[EFX $\fimplies$ PROPavg]
\label{thm:impl:efx-to-propavg}
Let $([n], [m], (v_i)_{i=1}^n, w)$ be a fair division instance
where $v_i$ is additive for some agent $i$.
If an allocation $A$ is EFX-fair to agent $i$, then it is also PROPavg-fair to agent $i$
if at least one of these conditions holds:
\begin{tightenum}
\item Agents have equal entitlements and $v_i(A_i) \ge 0$.
\item $n=2$.
\end{tightenum}
\end{lemma}
\begin{proof}
Suppose agent $i$ is EFX-satisfied but not PROPavg-satisfied by allocation $A$.
Then for some agent $\jhat \in [n] \setminus \{i\}$, we have
\[ \frac{v_i(A_i)}{w_i} < v_i([m]) < \frac{v_i(A_{\jhat})}{w_{\jhat}}. \]

Since $i$ doesn't EFX-envy $\jhat$, for all $S \subseteq A_i$ such that
$v_i(S \mid A_i \setminus S) < 0$, we get
\[ \frac{v_i(A_i \setminus S)}{w_i} \ge \frac{v_i(A_{\jhat})}{w_{\jhat}} > v_i([m]). \]
Hence, condition C2 of PROPavg is satisfied.
Since condition C4 of PROPavg is not satisfied, we get that
$T \neq \emptyset$ and $v_i(A_i) + \Sum(T)/|T| \le w_iv_i([m])$.
%
For all $\tau_j \in T$, we get
\begin{align}
\frac{v_i(A_i)}{w_i} &\ge \frac{\max(\{v_i(A_j \setminus S): S \subseteq A_j \textrm{ and } v_i(S \mid A_i) > 0\})}{w_j}
    \tag{$i$ doesn't EFX-envy $j$}
\\ &= \frac{v_i(A_j)}{w_j} - \frac{\min(\{v_i(S): S \subseteq A_j \textrm{ and } v_i(S) > 0\})}{w_j}
    \notag
\\ &= \frac{v_i(A_j)}{w_j} - \frac{\tau_j}{w_j}
    \label{eqn:impl:efx-to-propavg:1}
\end{align}

\textbf{Special case 1}: agents have equal entitlements and $v_i(A_i) \ge 0$.

Let $J \defeq \{j \in [n] \setminus \{i\}: \tau_j > 0\}$.
For all $j \in J$, using \eqref{eqn:impl:efx-to-propavg:1}, we get
$v_i(A_i) + \tau_j \ge v_i(A_j)$.

For all $j \in [n] \setminus (J \cup \{i\})$, we have $\tau_j = 0$, so $v_i(A_j) \le 0$. Hence,
\[ \frac{v_i(A_i)}{w_i} \ge 0 \ge \frac{v_i(A_j)}{w_j}. \]
Define $\tau_i \defeq 0$. Then,
\[ \frac{v([m])}{n} = \frac{1}{n}\sum_{j=1}^n v_i(A_j)
    \le \frac{1}{n}\sum_{j=1}^n (v_i(A_i) + \tau_j)
    = v_i(A_i) + \frac{\Sum(T)}{n} < v_i(A_i) + \frac{\Sum(T)}{|T|}. \]
Hence, condition 1 of PROPavg is satisfied, which is a contradiction.
Hence, it's impossible for $i$ to be EFX-satisfied by $A$ but not PROPavg-satisfied by $A$.

\textbf{Special case 2}: $n=2$

Let the two agents be $i$ and $j$. Since $T \neq \emptyset$, we have $T = \{\tau_j\}$, where $\tau_j > 0$.
Since $i$ is not PROPavg-satisfied by $A$, we get
\begin{align}
& v_i(A_i) + \tau_j \le w_iv_i([m]) = w_iv_i(A_i) + w_iv_i(A_j)
\\ &\implies w_jv_i(A_i) + \tau_j \le w_iv_i(A_j)
\\ &\implies \frac{v_i(A_i)}{w_i} \le \frac{v_i(A_j)}{w_j} - \frac{\tau_j}{w_iw_j}
    \le \frac{v_i(A_j)}{w_j} - \frac{\tau_j}{w_j}.
\end{align}
Combining this with equation \eqref{eqn:impl:efx-to-propavg:1} gives us $\tau_j = 0$, which is a contradiction.
Hence, it's impossible for $i$ to be EFX-satisfied by $A$ but not PROPavg-satisfied by $A$.
\end{proof}

\subsection{MMS vs EFX}
\label{sec:impls-extra:mms-vs-efx}

We prove some results connecting MMS, EFX, and related notions,
using techniques from \cite{plaut2020almost,caragiannis2023new}.

\begin{lemma}
\label{thm:mms-and-all-envy}
For a fair division instance $([n], [m], (v_i)_{i=1}^n, w)$,
suppose an allocation $A$ is WMMS-fair to agent $i$
and $i$ envies every other agent.
Then $A$ is also EFX-fair to agent $i$ if at least one of these conditions hold:
\begin{tightenum}
\item The items are goods to agent $i$ (i.e., $v_i(g \mid R) \ge 0$
    for all $R \subseteq [m]$ and $g \in [m] \setminus R$).
\item $v_i$ is additive and $w_i \le w_j$ for all $j \in [n] \setminus \{i\}$.
\end{tightenum}
\end{lemma}
\begin{proof}
Suppose agent $i$ is not EFX-satisfied by $A$, i.e., she EFX-envies some agent $j$.
Then $\exists S \subseteq A_i \cup A_j$ where either

\begin{tightenum}
\item $S \subseteq A_j$, $v_i(S \mid A_i) > 0$, and
    $\displaystyle \frac{v_i(A_i)}{w_i} < \frac{v_i(A_j \setminus S)}{w_j}$.
\item $S \subseteq A_i$, $v_i(S \mid A_i \setminus S) < 0$, and
    $\displaystyle \frac{v_i(A_i \setminus S)}{w_i} < \frac{v_i(A_j)}{w_j}$.
\end{tightenum}

If all items are goods, case 2 doesn't occur.

\textbf{Case 1}: $S \subseteq A_j$

Let $B$ be the allocation obtained by transferring $S$ from $A_j$ to $A_i$.
Formally, let $B_i \defeq A_i \cup S$, $B_j \defeq A_j \setminus S$,
and $B_k \defeq A_k$ for all $k \in [n] \setminus \{i, j\}$. Then
\[ \frac{v_i(B_i)}{w_i} = \frac{v_i(A_i) + v_i(S \mid A_i)}{w_i} > \frac{v_i(A_i)}{w_i}, \]
\[ \frac{v_i(B_j)}{w_j} = \frac{v_i(A_j \setminus S)}{w_j} > \frac{v_i(A_i)}{w_i}, \]
and for any $k \in [n] \setminus \{i, j\}$, we get
\[ \frac{v_i(B_k)}{w_k} = \frac{v_i(A_k)}{w_k} > \frac{v_i(A_i)}{w_i}. \]
Hence,
\[ \min_{k=1}^n \frac{v_i(B_k)}{w_k} > \frac{v_i(A_i)}{w_i} \ge \frac{\WMMS_i}{w_i}, \]
which is a contradiction.

\textbf{Case 2}: $S \subseteq A_i$

Let $v_i$ be additive and $w_i \le w_j$.
Let $B$ be the allocation obtained by transferring $S$ from $A_i$ to $A_j$. Formally,
let $B_i \defeq A_i \setminus S$, $B_j \defeq A_j \cup S$,
and $B_k \defeq A_k$ for all $k \in [n] \setminus \{i, j\}$. Then
\[ \frac{v_i(B_i)}{w_i} = \frac{v_i(A_i) - v_i(S)}{w_i} > \frac{v_i(A_i)}{w_i}, \]
\[ \frac{v_i(B_j)}{w_j} = \frac{v_i(A_j) + v_i(S)}{w_j} > \frac{v_i(A_i \setminus S)}{w_i} - \frac{(-v_i(S))}{w_j} \ge \frac{v_i(A_i)}{w_i}, \]
and for any $k \in [n] \setminus \{i, j\}$, we get
\[ \frac{v_i(B_k)}{w_k} = \frac{v_i(A_k)}{w_k} > \frac{v_i(A_i)}{w_i}. \]
Hence,
\[ \min_{k=1}^n \frac{v_i(B_k)}{w_k} > \frac{v_i(A_i)}{w_i} \ge \frac{\WMMS_i}{w_i}, \]
which is a contradiction.
\end{proof}

\begin{lemma}[MMS $\fimplies$ EFX for $n=2$]
\label{thm:impl:mms-to-efx-n2}
For a fair division instance $([2], [m], (v_i)_{i=1}^2, w)$,
suppose an allocation $A$ is WMMS-fair to agent 1.
Then $A$ is also EFX-fair to agent 1 if at least one of these conditions hold:
\begin{tightenum}
\item The items are goods to agent 1.
\item $v_1$ is additive and $w_1 \le w_2$.
\end{tightenum}
\end{lemma}
\begin{proof}
If agent 1 doesn't envy agent 2, she is EFX-satisfied.
Otherwise, she is EFX-satisfied because of \cref{thm:mms-and-all-envy}.
\end{proof}

Theorem 2 in \cite{caragiannis2023new} states that an MMS allocation is also an EEFX allocation
(for additive valuations over goods and equal entitlements).
The proof can be easily adapted to non-additive valuations over goods and unequal entitlements.
For the sake of completeness, we give a proof below.

\begin{lemma}[MMS $\fimplies$ EEFX, \cite{caragiannis2023new}]
\label{thm:impl:mms-to-eefx}
For a fair division instance $([n], [m], (v_i)_{i=1}^n, w)$,
if all items are goods to agent $i$ and allocation $A$ is WMMS-fair to agent $i$,
then $A$ is also epistemic-EFX-fair to $i$.
\end{lemma}
\begin{proof}
In any allocation $X$, let
$E_X$ be the set of agents envied by $i$,
$S_X$ be the set of agents EFX-envied by $i$,
$W_X$ be the total number of items among the agents in $S_X$.

$E_X \defeq \{t \in [n] \setminus \{i\}: i \textrm{ envies } t \textrm{ in } X\}$,
$S_X \defeq \{t \in [n] \setminus \{i\}: i \textrm{ EFX-envies } t \textrm{ in } X\}$,
$W_X \defeq \sum_{t \in S_X} |X_t|$, and $\phi(X) \defeq (-|E_X|, W_X)$.

First, we show that for any allocation $X$ where $|E_X| \le n-2$ and $S_X \neq \emptyset$,
there exists a \emph{better} allocation $Y$, i.e, $Y_i = X_i$ and $\phi(Y) < \phi(X)$
(tuples are compared lexicographically).

Let $j \in S_X$ and $k \in [n] \setminus \{i\} \setminus E_X$.
Since $i$ EFX-envies $j$, $\exists S \subseteq X_j$ such that $v_i(S \mid X_i) > 0$ and
\[ \frac{v_1(X_1)}{w_1} < \frac{v_1(X_j \setminus S)}{w_j}. \]
Let $Y_k \defeq X_k \cup S$, $Y_j \defeq X_j \setminus S$,
and $Y_t \defeq X_t$ for all $t \in [n] \setminus \{j, k\}$.

Then $i$ envies $j$ in $Y$. Hence, $E_X \subseteq E_Y$.
If $k \in E_Y$, then $|E_Y| > |E_X|$, so $\phi(Y) < \phi(X)$.
If $k \not\in E_Y$, then $W_Y < W_X$, so $\phi(Y) < \phi(X)$.
Hence, $Y$ is better than $X$.

Set $X = A$. As long as $|E_X| < n-1$ and $S_X \neq \emptyset$,
keep modifying $X$ as per Lemma 1.
This process will eventually end, since $\phi(X)$ keeps reducing,
and there are a finite number of different values $\phi(X)$ can take.
Let $B$ be the final allocation thus obtained.
Then $B_i = A_i$, and $|E_B| = n-1$ or $S_B = \emptyset$.

By \cref{thm:mms-and-all-envy},
$|E_B| = n-1$ implies $S_B = \emptyset$.
Hence, $B$ is agent $i$'s EEFX-certificate for $A$.
\end{proof}

\begin{definition}
\label{defn:leximin}
For any set $S \subseteq M$, let $\Pi_n(S)$ denote the set of all $n$-partitions of $S$.
For any sequence $X = (x_i)_{i=1}^n$ of real numbers, define $\sorted(X)$ to be a permutation of $X$
where entries occur in non-decreasing order.
For any two sequences $X = (x_i)_{i=1}^n$ and $Y = (y_i)_{i=1}^n$, we say that $X \le Y$ if
$\exists i \in [n]$ such that $x_i \le y_i$ and $x_j = y_j$ for all $j \in [i-1]$.
(Note that this relation $\le$ over sequences is a total ordering.)

We say that $P \in \Pi_n([m])$ is a leximin $n$-partition of a function $f: 2^{[m]} \to \mathbb{R}$ if
\[ P \in \argmax_{X \in \Pi_n([m])} \sorted\left((f(X_j))_{j=1}^n\right). \]
\end{definition}

It is easy to see that if $P$ is a leximin $n$-partition of $f: 2^{[m]} \to \mathbb{R}$,
then $\min_{j=1}^n f(P_j) = \MMS_f^n([m])$.

\begin{lemma}[MMS $\fimplies$ MXS]
\label{thm:impl:mms-to-mxs}
Let $\Ical \defeq ([n], [m], (v_i)_{i=1}^n, \eqEnt)$ be a fair division instance.
Let allocation $A$ be MMS-fair to agent $i$.
Then $A$ is also MXS-fair to $i$ if at least one of these conditions hold:
\begin{tightenum}
\item All items are goods to agent $i$.
\item $v_i$ is additive.
\end{tightenum}
\end{lemma}
\begin{proof}
We will show that agent $i$'s leximin $n$-partition is her MXS-certificate for $A$,
which would prove that $A$ is MXS-fair to agent $i$.

Without loss of generality, assume $i = 1$.
Let $P$ be a leximin $n$-partition of $v_1$ (c.f.~\cref{defn:leximin})
such that $v_1(P_1) \le v_1(P_2) \le \ldots \le v_1(P_n)$.
Then $v_1(P_1) = \MMS_{v_1}^n([m]) \le v_1(A_1)$.

For any agent $j \ge 2$, the allocation $(P_1, P_j)$ is leximin for
the instance $\Icalhat \defeq (\{i, j\}, P_1 \cup P_j, (v_1, v_j), (1/2, 1/2))$.
Hence, agent 1 is MMS-satisfied by $(P_1, P_j)$ in $\Icalhat$.
Since either all items are goods to agent 1 or $v_1$ is additive,
by \cref{thm:impl:mms-to-efx-n2}, agent 1 is EFX-satisfied by $(P_1, P_j)$ in $\Icalhat$.

On considering all values of $j$, we get that agent 1 is EFX-satisfied by $P$ in $\Ical$.
Hence, $P$ is agent 1's MXS-certificate for $A$.
\end{proof}

\subsection{Among PROP, APS, MMS}

We prove implications among share-based notions like PROP, APS, and MMS.

\begin{lemma}[PROP $\fimplies$ APS, Proposition 4 of \cite{babaioff2023fair}]
\label{thm:impl:prop-to-aps}
For any fair division instance
\\ $([n], [m], (v_i)_{i=1}^n, w)$,
$\APS_i \le w_iv_i([m])$ for agent $i$ if $v_i$ is additive.
\end{lemma}
\begin{proof}
Set the price $p(g)$ of each item $g$ to $v_i(g)$. Then
\[ \APS_i \le \max_{S \subseteq [m]: p(S) \le w_ip([m])} v_i(S)
= \max_{S \subseteq [m]: v_i(S) \le w_iv_i([m])} v_i(S) \le w_iv_i([m]).
\qedhere \]
\end{proof}

\begin{lemma}[PROP $\fimplies$ WMMS]
\label{thm:impl:prop-to-wmms}
For any fair division instance $([n], [m], (v_i)_{i=1}^n, w)$,
if $v_i$ is superadditive, then $\WMMS_i \le w_iv_i([m])$.
\end{lemma}
\begin{proof}
Let $P$ be agent $i$'s WMMS partition. Then
\[ v_i([m]) \ge \sum_{j=1}^n v_i(P_j) = \sum_{j=1}^n w_j\left(\frac{v_i(P_j)}{w_j}\right)
    \ge \sum_{j=1}^n w_j\frac{\WMMS_i}{w_i} = \frac{\WMMS_i}{w_i}. \qedhere \]
\end{proof}

\begin{lemma}
\label{thm:prefix-sum-bound}
Let $a_1 \le a_2 \le \ldots \le a_n$ be $n$ real numbers.
Let $s_k \defeq \sum_{i=1}^k a_i$ for any $0 \le k \le n$.
Then $s_k \le (k/n)s_n$.
\end{lemma}
\begin{proof}
$\displaystyle s_n = s_k + \sum_{i=k+1}^n a_i \ge s_k + (n-k)a_k \ge s_k + (n-k)\frac{s_k}{k}
= \frac{n}{k}s_k$.
\end{proof}

\begin{lemma}[PROP $\fimplies$ pessShare]
\label{thm:impl:prop-to-pessShare}
For any fair division instance $([n], [m], (v_i)_{i=1}^n, w)$,
if $v_i$ is superadditive, then $\pessShare_i \le w_iv_i([m])$.
\end{lemma}
\begin{proof}
Let $P$ be agent $i$'s $\ell$-out-of-$d$-partition.
Assume \wLoG{} that $v_i(P_1) \le v_i(P_2) \le \ldots \le v_i(P_d)$.
Then by \cref{thm:prefix-sum-bound} and superadditivity of $v_i$, we get
\[ \loodM_i = \sum_{j=1}^{\ell} v_i(P_j)
    \le \frac{\ell}{d}\sum_{j=1}^d v_i(P_j) \le \frac{\ell}{d}v_i([m]). \]
Hence,
\[ \pessShare_i \defeq \sup_{1 \le \ell \le d: \ell/d \le w_i} \loodM_i
    \le \sup_{1 \le \ell \le d: \ell/d \le w_i} (\ell/d)v_i([m]) \le w_iv_i([m]).
    \qedhere \]
\end{proof}

\begin{lemma}[APS $\fimplies$ pessShare]
\label{thm:impl:aps-to-pess}
For any fair division instance $([n], [m], (v_i)_{i=1}^n, w)$,
$\APS_i \ge \pessShare_i$ for every agent $i$.
\end{lemma}
\begin{proof}
Proposition 2 in \cite{babaioff2023fair} proves this for goods,
but their proof works for mixed manna too.
\end{proof}

\begin{lemma}
\label{thm:wmms-vs-knapsack}
For a fair division instance $([n], [m], (v_i)_{i=1}^n, w)$ and any agent $i$, define
\[ \beta_i \defeq \max_{j=1}^n \max_{\substack{S \subseteq [m]:\\v_i(S) \le w_jv_i([m])}} \frac{v_i(S)}{w_j}. \]
If $v_i$ is superadditive, then $\WMMS_i \le w_i\beta_i$.
If $v_i$ is additive and $n = 2$, then $\WMMS_i = w_i\beta_i$.
\end{lemma}
\begin{proof}
For any $j \in [n]$, define
\[ S_j \defeq \argmax_{\substack{S \subseteq [m]:\\v_i(S) \le w_jv_i([m])}} \frac{v_i(S)}{w_j}. \]
Then $\beta_i \defeq \max_{j=1}^n v_i(S_j)/w_j$.

Let $\Pi_n([m])$ be the set of all $n$-partitions of $[m]$. Let
\begin{align*}
P &\defeq \argmax_{P \in \Pi_n([m])} \min_{j=1}^n \frac{v_i(P_j)}{w_j},
& k &\defeq \argmin_{j=1}^n \frac{v_i(P_j)}{w_j}.
\end{align*}
Then $\WMMS_i/w_i \defeq v_i(P_k)/w_k$.

By \cref{thm:impl:prop-to-wmms}, we get $v_i(P_k)/w_k \le v_i([m])$. Hence,
\[ \frac{\WMMS_i}{w_i} = \frac{v_i(P_k)}{w_k} \le \frac{v_i(S_k)}{w_k}
    \le \max_{j=1}^n \frac{v_i(S_j)}{w_j} = \beta_i. \]

Now let $v_i$ be additive and $n = 2$.
For any $j \in [2]$, let $Q^{(j)}$ be an allocation where $Q^{(j)}_j \defeq S_j$
and $Q^{(j)}_{3-j} \defeq [m] \setminus S_j$.
Then $v_i(Q^{(j)}_j) = v_i(S_j) \le w_jv_i([m])$ and
\[ v_i(Q^{(j)}_{3-j}) = v_i([m]) - v_i(S_j) \ge v_i([m]) - w_jv_i([m]) = w_{3-j}v_i([m]). \]
Hence,
\[ \frac{v_i(Q^{(j)}_j)}{w_j} \le v_i([m]) \le \frac{v_i(Q^{(j)}_{3-j})}{w_{3-j}}. \]
Therefore,
\begin{align*}
\frac{\WMMS_i}{w_i} &= \min\left(\frac{v_i(P_1)}{w_1}, \frac{v_i(P_2)}{w_2}\right)
\\ &\ge \max_{j=1}^n \min\left(\frac{v_i(Q^{(j)}_j)}{w_j}, \frac{v_i(Q^{(j)}_{3-j})}{w_{3-j}}\right)
    \tag{by definition of $P$}
\\ &= \max_{j=1}^n \frac{v_i(S_j)}{w_j} = \beta_i.
\end{align*}
Hence, $\WMMS_i/w_i \ge \beta_i$.
\end{proof}

\begin{lemma}[WMMS $\fimplies$ APS for $n=2$]
\label{thm:impl:mms-to-aps-n2}
For a fair division instance $([2], [m], (v_i)_{i=1}^2, w)$,
if $v_i$ is additive for some $i$, then $\APS_i \le \WMMS_i$.
Moreover, when entitlements are equal, we get $\APS_i = \MMS_i$.
\end{lemma}
\begin{proof}
By \cref{thm:wmms-vs-knapsack}, we get
\[ \frac{\WMMS_i}{w_i} = \beta_i \defeq \max_{j=1}^2 \frac{v_i(S_j)}{w_j},
\quad\textrm{where}\quad
S_j \defeq \argmax_{\substack{S \subseteq [m]:\\ v_i(S) \le w_jv_i([m])}} v_i(S). \]
On setting $p = v_i$ in the primal definition of APS (\cref{defn:aps}), we get
\[ \APS_i \le \max_{\substack{S \subseteq [m]:\\ p(S) \le w_ip([m])}} v_i(S)
    = v_i(S_i) \le w_i\beta_i = \WMMS_i. \]
When entitlements are equal, $\APS_i \ge \pessShare_i \ge \MMS_i$
by \cref{thm:impl:aps-to-pess}.
\end{proof}

\subsection{Additive Triboolean Valuations}
\label{sec:impls-extra:tribool}

\begin{table*}[htb]
\centering
\caption[Tribool implications]{Implications among fairness notions when
valuations are additive and marginals are triboolean,
i.e, they belong to the set $\{-1, 0, 1\}$.}
\label{table:impls-tribool}
\small
\begin{tabular}{cccccccc}
\toprule & \footnotesize valuations & \footnotesize marginals & \footnotesize $n=2$ & \footnotesize entitlements & &
\\\midrule EF1 $\fimplies$ EFX & additive & $\{-1, 0, 1\}$
    & -- & -- & -- & trivial
\\[\defaultaddspace] EEF1 $\fimplies$ EEFX & additive & $\{-1, 0, 1\}$
    & -- & -- & -- & trivial
\\[\defaultaddspace] M1S $\fimplies$ MXS & additive & $\{-1, 0, 1\}$
    & -- & -- & -- & trivial
\\[\defaultaddspace] PROP1 $\fimplies$ PROPx & -- & $\{-1, 0, 1\}$
    & -- & -- & \cref{thm:impl:prop1-to-propx-tribool} & \textbf{new}
%
\\\midrule PROP $\fimplies$ EEF & additive & $\{-1, 0, 1\}$
    & -- & equal & \cref{thm:impl:tribool:prop} & \textbf{new}
\\[\defaultaddspace] APS $\fimplies$ PROPx & additive & $\{-1, 0, 1\}$
    & -- & -- & \cref{thm:impl:tribool:prop1,thm:impl:tribool:aps} & \textbf{new}
\\[\defaultaddspace] PROP1 $\fimplies$ APS & additive & $\{-1, 0, 1\}$
    & -- & -- & \cref{thm:impl:tribool:prop1,thm:impl:tribool:aps} & \textbf{new}
\\[\defaultaddspace] M1S $\fimplies$ APS & additive & $\{-1, 0, 1\}$
    & -- & equal & \cref{thm:impl:tribool:aps,thm:impl:tribool:m1s} & \textbf{new}
\\[\defaultaddspace] MMS $\fimplies$ EEFX & additive & $\{-1, 0, 1\}$
    & -- & equal & \cref{thm:impl:tribool:mms-to-eefx} & \textbf{new}
%
\\\midrule EF1 $\fimplies$ GAPS & additive & $\{-1, 0, 1\}$
    & -- & equal & \cref{thm:impl:tribool:ef1-gaps} & \textbf{new}
\\[\defaultaddspace] EF1 $\fimplies$ GAPS & additive & $\{-1, 0, 1\}$
    & $n=2$ & -- & \cref{thm:impl:tribool:ef1-gaps} & \textbf{new}
\\[\defaultaddspace] EF1 $\fimplies$ GAPS & additive & $\{-1, 0\}$
    & -- & -- & \cref{thm:impl:tribool:ef1-gaps} & \textbf{new}
\\ \bottomrule
\end{tabular}
\end{table*}

\begin{lemma}[PROP1 $\fimplies$ PROPx]
\label{thm:impl:prop1-to-propx-tribool}
Consider a fair division instance $([n], [m], (v_i)_{i=1}^n, w)$ where for some agent $i$,
marginals are triboolean, i.e., $v_i(t \mid S) \in \{-1, 0, 1\}$
for all $S \subseteq [m]$ and $t \in [m] \setminus S$ for some agent $i$.
Then if an allocation $A$ is PROP1-fair to agent $i$, then it is also PROPx-fair to $i$.
\end{lemma}
\begin{proof}
If $v_i(A_i) \ge w_iv_i([m])$, then $A$ is PROPx.
Otherwise, we get $v_i(A_i) + 1 > w_iv_i([m])$ since $A$ is PROP1
and marginals are triboolean.
Removing any positive-disutility subset from $A_i$
or adding any positive-utility subset to $A_i$ will increase its value by at least 1.
Hence, $A$ is PROPx.
\end{proof}

\begin{lemma}
\label{thm:tribool-rr}
For an additive function $f: 2^{[m]} \to \{-1, 0, 1\}$,
there exists an $n$-partition $P$ such that $|f(P_i) - f(P_j)| \le 1$
for all $i, j \in [n]$ and
and $\floor{f([m])/n} \le f(P_i) \le \ceil{f([m])/n}$ for all $i \in [n]$.
\end{lemma}
\begin{proof}
Partition $[m]$ into
goods $M_+ \defeq \{g \in [m]: v_i(g) > 0\}$,
chores $M_- \defeq \{c \in [m]: v_i(c) < 0\}$,
and neutral items $M_0 \defeq \{t \in [m]: v_i(t) = 0\}$.
Fuse items $M_0$, $\min(|M_+|, |M_-|)$ goods, and $\min(|M_+|, |M_-|)$ chores
into a single item $h$.
Then we are left with only goods and a neutral item,
or only chores and a neutral item.
Using round-robin, one can allocate items such that
any two bundles differ by at most one item.
\end{proof}

\begin{lemma}
\label{thm:impl:tribool:prop}
Consider a fair division instance $([n], [m], (v_i)_{i=1}^n, \eqEnt)$
where $v_i$ is additive and $v_i(t) \in \{-1, 0, 1\}$ for all $t \in [m]$ for some agent $i$.
If an allocation $A$ is PROP-fair to $i$, then it is also epistemic-EF-fair to $i$.
\end{lemma}
\begin{proof}
Since $A$ is PROP-fair to $i$, and $v_i(S) \in \mathbb{Z}$ for all $S \subseteq [m]$,
we get $v_i(A_i) \ge \ceil{v_i([m])/n}$.

Construct an allocation $B$ where $B_i = A_i$, and items $[m] \setminus A_i$
are allocated among agents $[n] \setminus \{i\}$ using \cref{thm:tribool-rr} with $f = v_i$.
We will show that $B$ is agent $i$'s epistemic-EF-certificate for $A$.

In $B$, for each agent $j \in [n] \setminus \{i\}$, we have
\[ v_i(B_j) \le \bigceil{\frac{v_i([m] \setminus A_i)}{n-1}}
    \le \bigceil{\frac{v_i([m]) - v_i([m])/n}{n-1}}
    \le \bigceil{\frac{v_i([m])}{n}} \le v_i(A_i). \]
Hence, $i$ doesn't envy anyone in $B$.
Hence, $B$ is agent $i$'s epistemic-EF-certificate for $A$.
\end{proof}

\begin{lemma}
\label{thm:impl:tribool:prop1}
Consider a fair division instance $([n], [m], (v_i)_{i=1}^n, w)$ where
$v_i$ is additive and $v_i(t) \in \{-1, 0, 1\}$ for all $t \in [m]$ for some agent $i$.
Then the following statements are equivalent:
\begin{tightenum}
\item Allocation $A$ is PROP1-fair to $i$.
\item Allocation $A$ is PROPx-fair to $i$.
\item $v_i(A_i) \ge \floor{w_iv_i([m])}$.
\end{tightenum}
\end{lemma}
\begin{proof}
Partition $[m]$ into
goods $M_+ \defeq \{g \in [m]: v_i(g) > 0\}$,
chores $M_- \defeq \{c \in [m]: v_i(c) < 0\}$,
and neutral items $M_0 \defeq \{t \in [m]: v_i(t) = 0\}$.

\textbf{Case 1}: $A_i$ has all goods and no chores.
\\ Then $v_i(A_i) \ge \max(0, v_i([m]))$.
If $v_i([m]) \ge 0$, then $v_i(A_i) \ge v_i([m]) \ge w_iv_i([m])$,
else $v_i(A_i) \ge 0 \ge w_iv_i([m])$.
Hence, $v_i(A_i) \ge w_iv_i([m])$ and $A$ is PROPx+PROP1.

\textbf{Case 2}: $A_i$ has a chore or some good is outside $A_i$.
\\ Then adding a good to $A_i$ or removing a chore from $A_i$
makes it's value more than $w_iv_i([m])$ iff $v_i(A_i) \ge \floor{w_iv_i([m])}$.
Hence, $A$ is PROP1 iff $A$ is PROPx iff $v_i(A_i) \ge \floor{w_iv_i([m])}$.
\end{proof}

\begin{lemma}
\label{thm:impl:tribool:aps}
Consider a fair division instance $([n], [m], (v_i)_{i=1}^n, w)$ where
$v_i$ is additive and $v_i(t) \in \{-1, 0, 1\}$ for all $t \in [m]$ for some agent $i$.
Then $\APS_i = \floor{w_iv_i([m])}$.
\end{lemma}
\begin{proof}
Partition $[m]$ into goods $M_+ \defeq \{g \in [m]: v_i(g) > 0\}$,
chores $M_- \defeq \{c \in [m]: v_i(c) < 0\}$, and
neutral items $M_0 \defeq \{t \in [m]: v_i(t) = 0\}$.

First, set $p(t) = v_i(t)$ for each $t \in [m]$ to get $\APS_i \le w_im_i$.
Since the APS is the value of some bundle, and bundle values can only be integers,
we get $\APS_i \le \floor{w_im_i}$.

Pick an arbitrary price vector $p \in \Delta_m$.
We will construct a set $S$ such that $p(S) \le w_ip([m])$ and $v_i(S) \ge w_iv_i([m])$.
Fuse items $M_0$, $\min(|M_+|, |M_-|)$ goods, and $\min(|M_+|, |M_-|)$ chores
into a single item $h$.
Let $M'_+$ and $M'_-$ be the remaining goods and chores, respectively.
Then $M'_+ = \emptyset$ or $M'_- = \emptyset$.
Using techniques from \cref{thm:aps-optimal-price},
we can assume \wLoG{} that $p_g \ge 0$ for all $g \in M'_+$,
$p_c \le 0$ for all $c \in M'_-$, and $p_h = 0$.

\textbf{Case 1}: $M'_- = \emptyset$.
\\ Let $m_i \defeq |M'_+| = v_i([m])$.
Let $S$ be the cheapest $\floor{w_im_i}$ items in $M'_+$.
Then using \cref{thm:prefix-sum-bound}, we get
\[ p(S) \le \frac{\floor{w_im_i}}{m_i}p(M_+) \le w_ip([m]). \]
Then $S$ is affordable and $v_i(S) = \floor{w_im_i}$.
Hence, $\APS_i \ge \floor{w_im_i}$.

\textbf{Case 2}: $M'_+ = \emptyset$.
\\ Let $m_i \defeq |M'_-| = -v_i([m])$.
Let $S$ be the cheapest $\ceil{w_im_i}$ items in $M'_- \cup \{h\}$.
Then using \cref{thm:prefix-sum-bound}, we get
\[ -p(S) \ge \frac{\ceil{w_im_i}}{m_i}(-p([m])) \ge w_i(-p([m])). \]
Then $S$ is affordable and $-v_i(S) \le \ceil{w_im_i} = -\floor{w_iv_i([m])}$.
Hence, $\APS_i \ge \floor{w_iv_i([m])}$.
\end{proof}

\begin{lemma}
\label{thm:impl:tribool:m1s}
Consider a fair division instance $([n], [m], (v_i)_{i=1}^n, \eqEnt)$
where $v_i$ is additive and $v_i(t) \in \{-1, 0, 1\}$ for all $t \in [m]$ for some agent $i$.
Then $\MMS_i = \mathrm{M1S}_i = \floor{v_i([m])/n}$.
\end{lemma}
\begin{proof}
Allocate items $[m]$ among agents $[n]$ using \cref{thm:tribool-rr} with $f = v_i$.
Then any two bundles differ by a value of at most one.
Hence, $\MMS_i = \floor{m'/n}$ and $\mathrm{M1S}_i \le \floor{m'/n}$,
where $m' \defeq v_i([m])$.

Let $X$ be an allocation where agent $i$ is EF1-satisfied.
Then any two bundles can differ by a value of at most one.
Hence, the smallest value $v_i(X_i)$ can have is $\floor{m'/n}$.
Hence, $\mathrm{M1S}_i \ge \floor{m'/n}$.
\end{proof}

\begin{lemma}
\label{thm:impl:tribool:ef1-gaps}
Consider a fair division instance $([n], [m], (v_i)_{i=1}^n, w)$ where
$v_i$ is additive and $v_i(t) \in \{-1, 0, 1\}$ for all $t \in [m]$ for some agent $i$.
If an allocation $A$ is EF1-fair to $i$, then it is also groupwise-APS-fair to $i$
if at least one of these conditions hold:
\begin{tightenum}
\item $n=2$
\item $w_i \le w_j$ for all $j \in [n] \setminus \{i\}$.
\item $v_i(c) \in \{0, -1\}$ for all $c \in [m]$.
\end{tightenum}
\end{lemma}
\begin{proof}
Consider any subset $S$ of agents.
On restricting $A$ to $S$, we get an allocation $B$ that is EF1-fair to $i$.
$B$ is also PROP1-fair to $i$ by \cref{thm:impl:eef1-to-prop1},
and APS-fair to $i$ by \cref{thm:impl:tribool:prop1,thm:impl:tribool:aps}.
Hence, $A$ is groupwise-APS-fair to $i$.
\end{proof}

\begin{remark}
\label{thm:ceil-floor}
For any $m \in \mathbb{Z}$ and $n \in \mathbb{Z}_{\ge 1}$, we get
\[ \bigfloor{\frac{m}{n}} = \bigceil{\frac{m+1}{n}} - 1. \]
\end{remark}

\begin{lemma}
\label{thm:impl:tribool:mms-to-eefx}
Consider a fair division instance $([n], [m], (v_i)_{i=1}^n, \eqEnt)$
where $v_i$ is additive and $v_i(t) \in \{-1, 0, 1\}$ for all $t \in [m]$ for some agent $i$.
If $v_i(A_i) \ge \floor{v_i([m])/n}$ for some allocation $A$,
then $A$ is epistemic-EFX-fair to $i$.
\end{lemma}
\begin{proof}
Construct an allocation $B$ where $B_i = A_i$, and items $[m] \setminus A_i$
are allocated among agents $[n] \setminus \{i\}$ using \cref{thm:tribool-rr} with $f = v_i$.
We will show that $B$ is agent $i$'s epistemic-EFX-certificate for $A$.

Let $k \defeq v_i([m])$.
Suppose $v_i(A_i) \ge \floor{k/n} = \ceil{(k+1)/n} - 1$ (c.f.~\cref{thm:ceil-floor}).
Then for any other agent $j \in [n] \setminus \{i\}$, we get
\[ v_i(B_j) \le \bigceil{\frac{v_i([m] \setminus A_i)}{n-1}}
    \le \bigceil{\frac{k - (k+1)/n + 1}{n-1}} = \bigceil{\frac{k+1}{n}}
    \le v_i(B_i) + 1. \]
If $B_j$ contains no goods and $B_i$ contains no chores,
then $v_i(B_i) \ge 0 \ge v_i(B_j)$, so $i$ doesn't envy $j$ in $B$.
Otherwise, transferring a good from $j$ to $i$ or a chore from $i$ to $j$ in $B$
eliminates $i$'s envy towards $j$.
Hence, $B$ is agent $i$'s epistemic-EFX-certificate for $A$.
\end{proof}

\subsection{Unit-Demand Valuations}
\label{sec:impls-extra:unit-demand}

We now study implications between fairness notions when agents have unit-demand valuations over goods.
A function $v: 2^{[m]} \to \mathbb{R}_{\ge 0}$ is said to be \emph{unit-demand} if $v(\emptyset) = 0$,
and for any non-empty set $X \subseteq M$, we have $v(X) = \max_{g \in X} v(\{g\})$.

Unit-demand functions capture scenarios where agents want just one good, e.g., when the goods are houses.
Unit-demand functions are known to be submodular and cancelable.
% TODO: prove it

\begin{table*}[htb]
\centering
\caption[Unit-demand implications]{Implications among fairness notions when
valuations are unit-demand, marginals are non-negative, and entitlements are equal.}
\label{table:impls-unit-demand}
\small
\begin{tabular}{ccHHHcH}
\toprule implication & \footnotesize valuations & \footnotesize marginals & \footnotesize $n=2$ & \footnotesize entitlements & proof &
\\\midrule MEFS $\fimplies$ EF
    & unit-demand & $\ge 0$ & -- & equal
    & \cref{thm:ud:ef} & trivial
\\[\defaultaddspace] MMS $\fimplies$ APS
    & unit-demand & $\ge 0$ & -- & equal
    & \cref{thm:ud:mms,thm:ud:aps} & \textbf{new}
\\[\defaultaddspace] PMMS $\fimplies$ PAPS
    & unit-demand & $\ge 0$ & -- & equal
    & \cref{thm:ud:mms,thm:ud:aps} & \textbf{new}
\\[\defaultaddspace] GMMS $\fimplies$ GAPS
    & unit-demand & $\ge 0$ & -- & equal
    & \cref{thm:ud:mms,thm:ud:aps} & \textbf{new}
\\[\defaultaddspace] M1S $\fimplies$ PROP1
    & unit-demand & $\ge 0$ & -- & equal
    & \cref{thm:ud:m1s,thm:ud:prop1} & \textbf{new}
\\[\defaultaddspace] M1S $\fimplies$ APS
    & unit-demand & $\ge 0$ & -- & equal
    & \cref{thm:ud:m1s,thm:ud:aps} & \textbf{new}
\\[\defaultaddspace] EF1 $\fimplies$ GAPS
    & unit-demand & $\ge 0$ & -- & equal
    & \cref{thm:ud:m1s,thm:ud:aps} & \textbf{new}
\\[\defaultaddspace] PPROP $\fimplies$ GPROP
    & unit-demand & $\ge 0$ & -- & equal
    & \cref{thm:impl:ud:pprop-to-gprop} & \textbf{new}
\\[\defaultaddspace] EF1 $\fimplies$ PROPavg
    & unit-demand & $\ge 0$ & -- & equal
    & \cref{thm:impl:ud:ef1-to-propavg} & \textbf{new}
\\ \bottomrule
\end{tabular}
\end{table*}

\begin{remark}
\label{thm:ud:ef}
Consider a fair division instance $([n], [m], (v_i)_{i=1}^n, \eqEnt)$ over goods
where $v_i$ is unit-demand for some agent $i$.
\WLoG, let $v_i(1) \ge v_i(2) \ge \ldots \ge v_i(m)$.
Then agent $i$ is envy-free in allocation $A$ iff $v_i(A_i) = v(1)$.
\end{remark}

\begin{lemma}
\label{thm:ud:m1s}
Consider a fair division instance $([n], [m], (v_i)_{i=1}^n, \eqEnt)$ over goods
where $v_i$ is unit-demand for some agent $i$.
\WLoG, let $v_i(1) \ge \ldots \ge v_i(m)$.
Then allocation $A$ is M1S-fair to agent $i$ iff $v_i(A_i) \ge v_i(n)$.
\end{lemma}
\begin{proof}
Suppose $A$ is M1S-fair to agent $i$. Let $B$ be agent $i$'s M1S-certificate for $A$.
Let $G \defeq \{g \in [m]: v_i(g) > v_i(B_i)\}$.
Then $B$ is EF1-fair to agent $i$, so every agent $j \in [n] \setminus \{i\}$
has at most one good from $G$, so $|G| \le n-1$.
If $v_i(B_i) < v_i(n)$, then $G \supseteq [n]$, which is a contradiction.
Hence, $v_i(A_i) \ge v_i(B_i) \ge v_i(n)$.

Suppose $v_i(A_i) \ge v_i(n)$. Let $B$ be an allocation where we give
exactly one good from $[n-1]$ to each agent in $[n] \setminus \{i\}$,
and allocate $[m] \setminus [n-1]$ to agent $i$.
Then $B$ is EF1-fair to $i$ and $v_i(A_i) \ge v_i(n) = v_i(B_i)$.
Hence, $B$ is an M1S-certificate for agent $i$ for allocation $A$.
Hence, $A$ is M1S-fair to agent $i$.
\end{proof}

\begin{lemma}
\label{thm:ud:mms}
Consider a fair division instance $([n], [m], (v_i)_{i=1}^n, \eqEnt)$ over goods
where $v_i$ is unit-demand for some agent $i$.
\WLoG, let $v_i(1) \ge \ldots \ge v_i(m)$.
Then agent $i$'s maximin share is $v_i(n)$.
\end{lemma}
\begin{proof}
In every $n$-partition $P$ of $[m]$, some bundle $P_k$ does not contain a good from $[n-1]$.
Then $v_i(P_k) \le v_i(n)$. Hence, $\MMS_i \le v_i(n)$.
Let $P$ be a partition where $P_j = \{j\}$ for $j \in [n-1]$ and $P_n = [m] \setminus [n-1]$.
Then $\min_{j=1}^n v_i(P_j) = v_i(n)$. Hence, $\MMS_i = v_i(n)$.
\end{proof}

\begin{lemma}
\label{thm:ud:aps}
Consider a fair division instance $([n], [m], (v_i)_{i=1}^n, \eqEnt)$ over goods
where $v_i$ is unit-demand for some agent $i$.
\WLoG, let $v_i(1) \ge \ldots \ge v_i(m)$.
Then agent $i$'s anyprice share is $v_i(n)$.
\end{lemma}
\begin{proof}
Set the price of the first $n-1$ goods to be $1/(n-1)$ each
and set the price of all the remaining goods to 0.
This tells us that $\APS_i \le v_i([m] \setminus [n-1]) = v_i(n)$.

For any non-negative price vector, some item from $[n]$ is affordable,
and gives value at least $v_i(n)$. Hence, $\APS_i \ge v_i(n)$.
\end{proof}

\begin{lemma}
\label{thm:ud:prop1}
Consider a fair division instance $([n], [m], (v_i)_{i=1}^n, \eqEnt)$ over goods
where $v_i$ is unit-demand for some agent $i$.
\WLoG, let $v_i(1) \ge \ldots \ge v_i(m)$.
Then every allocation is PROP1-fair to agent $i$.
\end{lemma}
\begin{proof}
Let $1 \in A_j$. If $j = i$, then $v_i(A_i) = v_i([m]) > v_i([m])/n$.
If $j \neq i$, then $v_i(A_i \cup \{1\}) = v_i([m]) > v_i([m])/n$.
\end{proof}

\begin{lemma}[PPROP $\fimplies$ GPROP]
\label{thm:impl:ud:pprop-to-gprop}
Consider a fair division instance $([n], [m], (v_i)_{i=1}^n, \eqEnt)$ over goods
where $v_i$ is unit-demand for some agent $i$.
\WLoG, let $v_i(1) \ge \ldots \ge v_i(m)$.
If an allocation $A$ is PPROP-fair to agent $i$, then it is also GPROP-fair to agent $i$.
\end{lemma}
\begin{proof}
Let $1 \in A_j$. $v_i(A_i) \ge v_i(A_i \cup A_j)/2 = v(1)/2$.
For any $S \subseteq [n]$ such that $i \in S$, we get
\[ \frac{1}{|S|}v_i\left(\bigcup_{j \in S} A_j\right)
\le \frac{v_i([m])}{2} = \frac{v(1)}{2} \le v_i(A_i). \]
Hence, $A$ is GPROP-fair to agent $i$.
\end{proof}

\begin{lemma}[EF1 $\fimplies$ PROPavg]
\label{thm:impl:ud:ef1-to-propavg}
Consider a fair division instance $([n], [m], (v_i)_{i=1}^n, \eqEnt)$ over goods
where $v_i$ is unit-demand for some agent $i$.
\WLoG, let $v_i(1) \ge \ldots \ge v_i(m)$.
If an allocation $A$ is EF1-fair to agent $i$, then it is also PROPavg-fair to agent $i$.
\end{lemma}
\begin{proof}
If $1 \in A_i$, then $A$ is PROP-fair. Now assume $1 \not\in A_i$.
Let $G \defeq \{g \in [m]: v_i(g) > v_i(A_i)\}$. Then $1 \in G$.
For each agent $j \in [n] \setminus \{i\}$, we have $|A_j \cap G| \le 1$
because $A$ is EF1-fair to agent $i$.
Let $J \defeq \{j \in [n] \setminus \{i\}: A_j \cap G \neq \emptyset\}$.

Recall the definition of $\tau_j$ from \cref{defn:propavg}.
If $A_j \cap G = \emptyset$, then $\tau_j = 0$.
If $A_j \cap G \neq \emptyset$, then $\tau_j = v_i(A_j) - v_i(A_i) > 0$.
%
\begin{align*}
& v_i(A_i) + \frac{1}{n-1}\sum_{j \in [n] \setminus \{i\}} \tau_j
= v_i(A_i) + \frac{1}{n-1}\sum_{j \in J} (v_i(A_j) - v_i(A_i))
\\ &= \left(1 - \frac{|J|}{n-1}\right)v_i(A_i) + \frac{1}{n-1}\sum_{g \in G} v_i(g)
\ge \frac{v_i(1)}{n-1} > \frac{v_i([m])}{n}.
\end{align*}
Hence, $A$ is PROPavg-fair to agent $i$.
\end{proof}
